\chapter{Background and Theoretical Foundations}
\label{ch:background}

\section{LoRaWAN Technology}
\label{sec:lorawan-tech}

\subsection{LoRa Physical Layer}
% CSS modulation, SF7–SF12, EU868 (863–870 MHz), bandwidth 125/250 kHz

\subsection{LoRaWAN Protocol Architecture}
% End device → gateway → network server → application server
% Class A (our devices): uplink-initiated, two Rx windows after each uplink
% OTAA vs ABP — we use OTAA

\subsection{Duty Cycle and Payload Constraints}
% EU868: 1% duty cycle per sub-band
% Payload by DR: DR0/SF12 = 51 bytes (our worst-case constraint)
% This shaped our 52-byte FL payload design: 4 header + 48 weight deltas
\begin{table}[H]
  \centering
  \begin{tabular}{cccc}
    \toprule
    \textbf{DR} & \textbf{SF} & \textbf{Max Payload (B)} & \textbf{Approx.\ Airtime (51\,B)} \\
    \midrule
    0 & 12 & 51  & $\approx$2.5\,s \\
    1 & 11 & 51  & $\approx$1.3\,s \\
    2 & 10 & 51  & $\approx$0.6\,s \\
    3 & 9  & 115 & $\approx$0.3\,s \\
    4 & 8  & 222 & $\approx$0.2\,s \\
    5 & 7  & 222 & $\approx$0.1\,s \\
    \bottomrule
  \end{tabular}
  \caption{LoRaWAN EU868 maximum payload size by data rate.}
  \label{tab:lorawan-payload}
\end{table}

\subsection{The Things Network}
% TTN v3 infrastructure used in both prior and current project
% Webhook integration: TTN → POST /uplink → fl_server.py
% Downlink limit: ~10 confirmed downlinks/day/device

\section{Indoor Radio Propagation and Path Loss}
\label{sec:path-loss-theory}

\subsection{Free-Space and Log-Distance Path Loss Models}
% Friis: PL_fs = 20·log10(d) + 20·log10(f) + 20·log10(4π/c)
% Log-distance: PL(d) = PL(d_0) + 10·n·log10(d/d_0) + X_σ
% n ≈ 2.0–3.5 for indoor office environments

\subsection{Expected Path Loss}
% exp_pl = TX_power − RSSI + antenna_corrections
% This is the target variable for our regression model

\subsection{Environmental Factors Affecting Indoor Propagation}
% Prior project finding: humidity, CO2, PM2.5, temperature, pressure
% all correlate with path loss
% Physical mechanisms: water absorption, occupancy effects, refraction

\section{Machine Learning for Path Loss Prediction}
\label{sec:ml-pathloss}

\subsection{Traditional Centralized Approaches}
% MLR, Random Forest, XGBoost — prior project achieved R² ≈ 0.93 with XGBoost

\subsection{Neural Network Approaches}
% Dense NNs for regression; knowledge distillation for edge deployment
% Motivation for Dense(8→8→1): smallest model capturing env+radio→PL relationship

\section{Federated Learning}
\label{sec:fl-theory}

\subsection{Fundamentals}
% McMahan et al. (2017): train locally, aggregate globally, no raw data sharing

\subsection{Federated Averaging Algorithm}
\begin{algorithm}[H]
  \caption{\gls{fedavg}~\cite{mcmahan2017communication}}
  \label{alg:fedavg}
  \KwIn{$K$ clients, model $w_0$, $T$ rounds, $E$ local epochs, lr $\eta$}
  \KwOut{Global model $w_T$}
  \For{round $t = 1 \ldots T$}{
    Server selects subset $S_t \subseteq \{1,\ldots,K\}$\;
    Server broadcasts $w_t$ to $S_t$\;
    \For{client $k \in S_t$ \textbf{in parallel}}{
      $w_k \leftarrow w_t$\;
      \For{epoch $e = 1 \ldots E$}{
        \For{batch $\mathcal{B} \subset \mathcal{D}_k$}{
          $w_k \leftarrow w_k - \eta \nabla \ell(w_k;\mathcal{B})$\;
        }
      }
      Send $w_k$ to server\;
    }
    $w_{t+1} \leftarrow \displaystyle\sum_{k \in S_t} \frac{n_k}{\sum_{j} n_j} w_k$\;
  }
\end{algorithm}

\subsection{Non-IID Data Challenges}
% Real deployments: each device sees different data; our case: ED0 (stairwell)
% has different conditions than ED3 (kitchen); causes slower convergence

\subsection{Communication Efficiency: Quantization and Compression}
% float32 → int8 (scale factor 127); critical for 51-byte LoRaWAN limit

\section{TinyML}
\label{sec:tinyml-theory}

\subsection{TensorFlow Lite Micro}
% Static memory via tensor arena; AllOpsResolver; no OS needed

\subsection{Model Quantization for Microcontrollers}
% Post-training int8 quantization with representative dataset calibration

\subsection{On-Device Training Challenges}
% TFLite Micro has NO backpropagation support
% Limited RAM prevents full gradient computation
% Current approach: simplified weight update proxy


% ══════════════════════════════════════════════════════════════════════════════
% CHAPTER 3: RELATED WORK
% ══════════════════════════════════════════════════════════════════════════════
